\subsection{The final dyadic residual branch and its exact least index}

\begin{lemma}[B-RD48: the marked level-eight square conditions]
\label{lem:B-RD48-models}
\status{Proved; complete independent proof, source and code audits passed.}
Let $E:y^2=x^3+A x+B$ be an elliptic curve over a field $K$ of
characteristic zero, let $\theta$ be a root of its cubic, and put
\[
 \mathcal D=-4A^3-27B^2,\qquad
 \gamma=-(3\theta^2+A),\qquad
 q_\theta(t)=t^2-\frac{2\theta}{3}t-
                         \frac{2\theta^2+A}{9}.          \tag{B-RD48.1}
\]
If an elliptic curve $E'/K$ has an isomorphism $E[8]\simeq E'[8]$
of Weil-pairing multiplier $r\in\{1,3,5,7\}$, its induced
marked level-four point has a coordinate $t\in\mathbf P^1(K)$.
For finite $t$, the following element is a square in $K(\theta)$:
\[
 \alpha_r q_\theta(t),\qquad
 \begin{array}{c|rrrr}
 r&1&5&3&7\\\hline
 \alpha_r&1&\mathcal D&\mathcal D/\gamma&\gamma.
 \end{array}                                             \tag{B-RD48.2}
\]
For the level-four family $E_t$ with $E_\infty=E$, let
\[
 \Phi_E(t)=\operatorname{Norm}_{K(\theta)/K}q_\theta(t)
                                                               \tag{B-RD48.3}
\]
when the cubic is irreducible. Then
\[
                  \Delta(E_t)=\Delta(E)\Phi_E(t)^4.      \tag{B-RD48.4}
\]
The curve $E'$ is $E_t$ for $r=1,5$, and the quadratic twist
$E_t^{\Delta(E)}$ for $r=3,7$, in the compatible families.
At $t=\infty$ these are respectively $E$ and $E^{\Delta(E)}$.
In particular the discriminant valuations of $E'$ and $E_t$
are congruent modulo three over any discretely valued $K$.
\end{lemma}
\begin{proof}
Use Chen's marked moduli models
\cite[Theorems~1.1--1.4, pp.2--3; Section~3, pp.5--7;
Sections~5--6, pp.10--16; Appendix, p.21]{Chen2014Mod8}.
Here we spell out both the field scope and the normalization.
The finite-level moduli problem records an actual isomorphism
of eight-torsion with prescribed Weil-pairing multiplier.
Its forgetful map records the induced four-torsion marking.
Chen constructs its covers of the level-four parameter line
by square roots and computes the twists by cocycles in the
finite group $\GL_2(\mathbf Z/8)$.
The calculation is performed on generators of that finite
group (Lemmas~5.1--5.5 and 6.1--6.3), so it can be restricted
to the image of $G_K$ for any characteristic-zero field $K$.
Equivalently, these are identities for the universal marked
elliptic curve over
$\mathbf Q[A,B,(4A^3+27B^2)^{-1}]$, followed by base change.
They do not require that $A,B$ be algebraic over $\mathbf Q$,
and they apply at $j=0$ and $1728$ as well. Nontrivial elliptic
automorphisms cause no extra quadratic ambiguity: a full
level-four marking has no nonidentity automorphism fixing it.

The level-four completeness and field scope also follow directly
from Fisher's Theorem~13.2, stated for arbitrary $K$, with its
proof and the discussion of the special $j$ values
\cite[Theorem~13.2 and proof, preprint pp.29--30;
discussion and reverse-four twist, p.31]{Fisher2010Hessian}.
In Chen's coordinate change of Section~3, put
$T(u,v)=uv(16u^4-v^4)$.
Its matrix is
$\left(\begin{smallmatrix}u&-T_v\\v&T_u\end{smallmatrix}\right)$,
whose determinant is $uT_u+vT_v=6T$.
The original elliptic discriminant is nonzero, so $T\ne0$.
Thus this coordinate change is invertible also at $j=0,1728$;
the exceptional-coordinate restriction of Fisher's
Corollary~13.4 is not being used.

For clarity, the equation for multiplier one, after comparing
coefficients of $1,\theta,\theta^2$, is exactly
\[
                  (a_0+a_1\theta+a_2\theta^2)^2
                                =q_\theta(t).
\]
The other three displayed models give the scaling factors in
(B-RD48.2). In particular the coefficient of $t^2$ in the
multiplier-three model is
\[
                    6A\theta^2-9B\theta+4A^2
                              =\mathcal D/\gamma,
\]
using $\theta^3+A\theta+B=0$; its remaining coefficients
are those of $\mathcal D q_\theta/\gamma$.
The multiplier-seven model similarly gives $\gamma q_\theta$.
This proves the necessary square conditions, including the
actual reversed markings. The quadratic twist for the
reversed level-four family is Chen's Section~3 and Lemma~3.1.

Here is an independent polynomial verification of the
discriminant normalization. Write $c=-A/27$, $d=-B/54$.
The norm in (B-RD48.3) is
\[
 \Phi_E(t)=t^6-15ct^4-40dt^3-45c^2t^2-24cdt+27c^3-32d^2.
                                                               \tag{B-RD48.5}
\]
Let $a_E(t),b_E(t)$ be the two polynomials in Chen's Appendix,
so that $E_t$ has coefficients $-27a_E(t),-54b_E(t)$.
Expansion gives the universal identity
\[
       a_E(t)^3-b_E(t)^2=(c^3-d^2)\Phi_E(t)^4,
\]
which is (B-RD48.4). All its coefficients, and the coefficient
comparisons above, are checked exactly in the retained certificate.
This is independently the degree-four instance of Fisher's
Hesse identity
\cite[Lemma~8.4, equation~(8.1), and Remark~8.6,
preprint pp.19--20]{Fisher2010Hessian}.
The leading terms of the two Appendix polynomials are
$ct^8,dt^{12}$, so the family at infinity is $E$ after the
coordinate scaling by $t^2$.

\emph{Source sign check.} The minus sign in (B-RD48.1) is
the one in Chen's Theorems~1.1--1.4 and agrees with the
Appendix discriminant identity just verified. The cusp-sum
formula in printed Lemma~3.2 has the opposite sign.
We do not import that sign into the computation. Replacing
$t$ by $-t$ changes no valuation used below, but the displayed
models and the certificate consistently use (B-RD48.1).
\end{proof}

\begin{proposition}[B-RD48: the residual cubic branch is excluded]
\label{prop:B-RD48}
\status{Proved; complete independent proof, source and code audits passed.}
Let $K$ be a complete discretely valued field of mixed
characteristic $(0,2)$ with perfect residue field and $v_K(2)=1$.
If two potentially supersingular elliptic curves have
isomorphic eight-torsion and common residual inertia image
$C_3$, their special-inertia characters agree on the elements
of odd order in a common finite special-inertia quotient.
\end{proposition}
\begin{proof}
Pass to the completed maximal unramified extension.
Its residue field is algebraically closed, every unit is a cube
by Hensel's lemma, and it contains $\zeta_3$.
The common two-torsion field is the unique tame cubic extension
\[
                   F=K(\alpha),\qquad \alpha^3=2,
                   \qquad v_F(2)=3.
\]
Suppose for a contradiction that the characters differ, and
name $E$ as the curve whose finite special-inertia image has
order prime to three. This is possible by
Lemma~\ref{lem:B-RD46-discriminant} and its character conclusion.

Complete the square and remove the quadratic term in a
Weierstrass equation. Its roots lie in $F$, have trace zero,
and form a three-cycle. Thus one root has the form
$a\alpha+b\alpha^2$ with $a,b\in K$.
The terms, when nonzero, have distinct valuations modulo three.
Direct multiplication of the three roots gives
\[
 A=-6ab,\qquad B=-2a^3-4b^3,\qquad
 \Delta=-1728(a^3-2b^3)^2.                             \tag{B-RD48.6}
\]
If $r_0=\min\{3v_K(a)+1,3v_K(b)+2\}$, the same formula gives
$v_K(\Delta)=4+2r_0$.
This valuation is divisible by three, by our choice of $E$.
Consequently the first term is the smaller one, $a\ne0$,
and $v_K(b)\ge v_K(a)$.
Twist both curves and their eight-torsion isomorphism by
$a^{-1}$. Such a common quadratic twist does not change the
characters on odd-order elements. The first curve becomes
\[
 E_u:\quad y^2=x^3-6u x-2-4u^3,
 \qquad u=b/a\in\mathcal O_K,
 \qquad \theta=\alpha+u\alpha^2.                       \tag{B-RD48.7}
\]
Here $v_F(\theta)=1$, $v_K(\Delta(E_u))=6$, and
$v_F(\mathcal D)=6$. Also
\[
             v_F(\gamma)=v_F(3\theta^2+A)=2,
\]
because the other term has valuation at least three.
The four scaling factors in (B-RD48.2) therefore have
valuations $0,6,4,2$, all even.

Let $t$ be the marked level-four coordinate of the second
curve supplied by Lemma~\ref{lem:B-RD48-models}.
The point at infinity gives $E_u$ or its discriminant twist,
which has the same odd-order character, so it is already
inconsistent with our supposition. Assume $t\in K$ and put
$k=v_K(t)$ when $t\ne0$.
The four terms of $q_\theta(t)$ in (B-RD48.1) have valuations
\[
                6k,\qquad4+3k,\qquad5,
                         \qquad3+3v_K(u),              \tag{B-RD48.8}
\]
with the last term omitted when $u=0$.
If $k\le0$, the first is the unique minimum, so
$v_F(q_\theta(t))=6k$.
If $k\ge1$, the unique minimum is three when $u$ is a unit,
and five when $u$ is not a unit. The same assertion holds
at $t=0$ after omitting its first two terms.
In the latter cases every
$\alpha_rq_\theta(t)$ has odd valuation in $F$ and hence
cannot be a square. This contradicts the actual level-eight
marking, for every possible multiplier $r$.

Thus $k\le0$. Since $F/K$ is totally ramified with residue
degree one,
\[
          v_K\!\left(\operatorname{Norm}_{F/K}q_\theta(t)\right)
                              =v_F(q_\theta(t))=6k.
\]
The exact discriminant identity gives
\[
                         v_K(\Delta(E_t))=6+24k,
\]
which is divisible by three. For a reversed multiplier the
additional discriminant twist changes this by a multiple of
six, hence preserves divisibility by three.
Lemma~\ref{lem:B-RD46-discriminant} says that the second
potentially supersingular curve has no cubic special inertia
either. This is the desired contradiction.
\end{proof}

\begin{theorem}[B-RD48: the sharp dyadic height-two index]
\label{thm:B-RD48}
\status{Proved; complete independent proof, source and code audits passed.}
In the notation and monodromy class of B.4,
\[
                            \nu_{\min}(2,1,2)=3.
\]
Thus sixteen-torsion determines the requested prime-to-two
special-inertia character of every permitted height-two group
over absolute ramification index one, while eight-torsion does not.
For the subclass consisting of two potentially supersingular
elliptic BT groups, eight-torsion already determines that character.
\end{theorem}
\begin{proof}
For two potentially supersingular elliptic curves, the residual
inertia images over the completed maximal unramified field are
only $1,C_2,C_3$. The trivial case follows from
Proposition~\ref{prop:B-RD46}, the order-two case from
Proposition~\ref{prop:B-RD47}, and the cubic case from the
preceding proposition. This proves the stated subclass result.

For arbitrary permitted height-two groups with isomorphic
sixteen-torsion, Proposition~\ref{prop:B-RD44} already excludes
every discrepancy except a comparison of two exact potentially
supersingular elliptic realizations. Apply the subclass result
to their eight-torsion. This proves the upper bound three in
the manuscript's indexing. Proposition~\ref{prop:B-RD43}
supplies the explicit ramification-one eight-torsion comparison
with unequal characters, proving the lower bound three.
No assertion about the least index for greater height or
greater absolute ramification follows from this small instance.
\end{proof}

\paragraph{Exact certificate and current status.}
The complete retained certificate
\path{code/verify_p2_residual_c3_gate.py}
checks the full universal discriminant polynomial, all needed
coefficient and scaling identities, the trace-zero normal form,
and $20\,301$ valuation cases together with $3\,721$ normalization
cases. It passed in $0.6555589939962374$ seconds.
The finite arithmetic is additional verification of the symbolic
proof; no finite sample replaces the parameter or moduli arguments.
All three assertions have passed complete independent proof,
source and code audits. The parent reran the complete certificate
in $0.8482220709993271$ seconds, with its audit recorded at
2026-11-19 14:30:51.582619 UTC. The obstructions/crystals agent
independently reread the original Chen and Fisher passages and
reran the full certificate in $0.5984987490010099$ seconds.
Both preserved the author output. The source-sign discrepancy
and its exact resolution remain recorded above; no failed
normalization has been silently discarded.
